\paragraph{Пример.}
Для запроса 1086224 «what causes boiling point to increase» («что вызывает повышение температуры кипения») в таблице~\ref{tab:annotation-example} сопоставлены исходный релевантный фрагмент, дополнительный правильный ответ и отрицательный кандидат. Описания являются сокращённым пересказом сохранённых фрагментов.

\begin{table}[t]
\centering
\caption{Конкретный пример доразметки. «Не размечен» означает отсутствие в исходных положительных qrels, а не обязательно нерелевантность.}
\label{tab:annotation-example}
\small
\begin{tabularx}{\textwidth}{@{}l L l l@{}}
\toprule
ID & Содержание фрагмента & Исходная & Ручная \\
\midrule
7112148 & Добавление соли повышает температуру кипения воды. & Релевантен & Релевантен \\
7399382 & Нелетучее растворённое вещество, например соль, повышает температуру кипения растворителя. & Не размечен & Релевантен \\
987600 & Приведена температура кипения ацетона. & Не размечен & Нерелевантен \\
\bottomrule
\end{tabularx}
\end{table}

Если система возвращает фрагмент 7399382, исходная неполная разметка не даёт ему положительного вклада. После дополнительной ручной оценки тот же фрагмент даёт положительный вклад без изменения ранжирования системы.
